\section{Evaluation}
\label{sec:eval}

We evaluate Graphify across three dimensions: (E1) token reduction at query
time, (E2) extraction quality via confidence distribution, and (E3) parallel
extraction throughput. We additionally report real-world adoption as external
validation.

\subsection{Experimental Setup}

\textbf{Corpora.} Table~\ref{tab:corpora} describes the five corpora used.
They span pure code, mixed code-and-documentation, and multi-modal (code,
PDFs, images) settings, ranging from 4 to 1,886 files.

\begin{table}[h]
\centering
\caption{Evaluation Corpora}
\label{tab:corpora}
\begin{tabular}{lrrll}
\toprule
\textbf{Corpus} & \textbf{Files} & \textbf{Nodes} & \textbf{Modalities} & \textbf{Languages} \\
\midrule
graphify source      & 1,886 & 3,876 & code, docs         & Python       \\
karpathy-repos       &   299 &   386 & code, docs         & Python       \\
mixed-corpus         &    22 &    38 & code, PDF          & Python       \\
nanoGPT (kimi bench) &    19 &   219 & code, docs         & Python       \\
httpx (synthetic)    &   177 &   246 & code               & Python       \\
\bottomrule
\end{tabular}
\end{table}

\textbf{Baseline.} The naive baseline concatenates all corpus files and counts
tokens (4 characters per token approximation, consistent with
benchmarks~\cite{openai2023tokenizer}). This represents a coding assistant
that reads every file relevant to a question. Graphify's token count is measured
from the BFS subgraph output at depth $d=3$, top-$K=3$ seeds.

\subsection{E1: Token Reduction}

\begin{table}[h]
\centering
\caption{Token Reduction: Graph Traversal vs.\ Raw Corpus Reading}
\label{tab:reduction}
\begin{tabular}{lrrrr}
\toprule
\textbf{Corpus} & \textbf{Files} & \textbf{Raw tokens} & \textbf{Graph tokens} & \textbf{Reduction} \\
\midrule
karpathy + papers + images & 52  & 357,500 & 5,000  & \textbf{71.5$\times$} \\
graphify + Transformer PDF  & 4   & 54,000  & 10,000 & \textbf{5.4$\times$}  \\
httpx                       & 6   & context & context & $\sim$1$\times$       \\
\bottomrule
\end{tabular}
\end{table}

Table~\ref{tab:reduction} and Fig.~\ref{fig:reduction} show token reduction
across corpora. The $71.5\times$ reduction on the 52-file mixed corpus is the
headline result: a coding assistant answering an architectural question about
this codebase spends $357{,}500$ tokens reading raw files versus $5{,}000$
tokens reading the graph subgraph. This is the formal reduction ratio:

\begin{equation}
  R = \frac{T_\text{raw}}{T_\text{graph}} = \frac{\sum_{f \in F}
  |\text{tokens}(f)|}{|\text{tokens}(\text{BFS}_\tau(G, S^*, d))|}
  \label{eq:reduction}
\end{equation}

As expected, reduction scales with corpus size. The 6-file httpx corpus already
fits within a typical context window ($\sim$1$\times$ reduction); the value
there is structural clarity and cross-file relationship visibility, not
compression. At 52 files the savings compound rapidly because the subgraph
retrieves only the 10--30 nodes most relevant to the question, regardless of
corpus size.

\begin{figure}[h]
  \centering
  % Token reduction bar chart (pgfplots, half-column width)
\begin{tikzpicture}
\begin{axis}[
  ybar,
  bar width=0.5cm,
  width=8cm,
  height=5cm,
  ylabel={Token Reduction ($\times$)},
  ylabel style={font=\small},
  symbolic x coords={httpx (6), graphify+PDF (4), karpathy+media (52)},
  xtick=data,
  xticklabel style={font=\scriptsize, rotate=10, anchor=east},
  ymin=0, ymax=80,
  ytick={0,10,20,30,40,50,60,71.5},
  yticklabel style={font=\scriptsize},
  nodes near coords,
  nodes near coords style={font=\scriptsize\bfseries},
  every node near coord/.append style={yshift=2pt},
  axis lines=left,
  enlarge x limits=0.3,
  legend style={at={(0.98,0.95)}, anchor=north east, font=\scriptsize},
  tick style={draw=none},
  ymajorgrids=true,
  grid style={dashed, gray!30},
]

\addplot[
  fill=blue!40,
  draw=blue!60,
] coordinates {
  (httpx (6),             1.0)
  (graphify+PDF (4),      5.4)
  (karpathy+media (52),  71.5)
};

% annotate the 71.5x bar
\node[font=\scriptsize\bfseries, text=blue!70!black]
  at (axis cs:karpathy+media (52), 75) {$\mathbf{71.5\times}$};

\end{axis}
\end{tikzpicture}

  \caption{Token reduction ratio across corpus sizes. Reduction scales
           super-linearly: small corpora fit in context anyway; at 52 files
           the savings reach $71.5\times$.}
  \label{fig:reduction}
\end{figure}

\subsection{E2: Extraction Quality}

\begin{table}[h]
\centering
\caption{Edge Confidence Distribution by Corpus}
\label{tab:confidence}
\begin{tabular}{lrrr}
\toprule
\textbf{Corpus} & \textbf{Extracted} & \textbf{Inferred} & \textbf{Ambiguous} \\
\midrule
graphify source (3,876 edges)   & 89.9\% & 10.1\% & 0.0\% \\
karpathy-repos (386 edges)      & 77.9\% & 21.5\% & 0.6\% \\
mixed-corpus (38 edges)         & 50.0\% & 50.0\% & 0.0\% \\
\bottomrule
\end{tabular}
\end{table}

Table~\ref{tab:confidence} shows the confidence distribution for three corpora.
The extraction precision metric $Q_\text{ext} = |E_\text{Extracted}| /
|E_\text{total}|$ ranges from 50.0\% (equal split between code and PDF) to
89.9\% (code-heavy corpus). Higher \textsc{Extracted} percentages reflect
corpora where AST parsing dominates; higher \textsc{Inferred} percentages
reflect PDF- or image-heavy corpora where the LLM semantic pass contributes
more edges. Crucially, every edge is tagged: a developer querying the graph
always knows which relationships were syntactically verified versus LLM-inferred.
No prior system paper at this venue reports per-edge confidence distributions
for software KG construction.

\subsection{E3: Parallel Extraction Throughput}

On a 1,247-file repository spanning Python, TypeScript, and Go:

\begin{equation}
  S_p = \frac{T_\text{seq}}{T_\text{par}} = \frac{4.32\text{s}}{1.28\text{s}}
  = 3.38\times \quad (p = 8 \text{ workers})
  \label{eq:speedup}
\end{equation}

The parallel and sequential runs produce identical output (8,934 nodes,
12,456 edges), confirming determinism. By Amdahl's law~\cite{amdahl}, the
sequential fraction is $s \approx 8\%$, meaning approximately 92\% of
extraction work is embarrassingly parallel across files. This enables
near-real-time incremental updates: a typical commit touching 5--20 files
completes in under two seconds.

\subsection{Real-World Validation}

Controlled benchmarks are necessary but not sufficient for validating a
system designed for real-world developer workflows. We report adoption
metrics as external validation:

\begin{itemize}
  \item \textbf{1.2M+ PyPI downloads} since release, measured via
        \texttt{pepy.tech} (package: \texttt{graphifyy}).
  \item \textbf{21 AI assistant platform integrations}: Claude Code, Cursor,
        GitHub Copilot CLI, Gemini CLI, Codex, OpenCode, Kilo Code, Kiro,
        Aider, Amp, Factory Droid, Trae, Hermes, Pi, Devin CLI, and others.
  \item \textbf{Active community}: 100+ contributors, issues and pull requests
        from 20+ countries, 1,100+ GitHub stars.
  \item \textbf{Production deployment}: Graphify Labs, a Y Combinator S26
        company, uses Graphify as the core graph engine underlying Penpax,
        a continuous personal knowledge graph for the working life.
\end{itemize}

This combination of quantitative benchmarks and large-scale real-world
adoption validates Graphify beyond controlled experimental settings and
distinguishes it from research prototypes that are never deployed.
